\section*{四、实验探讨与思考题}

\subsection*{1. 实验小结}

本次实验按课堂要求完成手动前后两次自动测量，以及一次覆盖 0.5—100.0 V 的完整手动测量。从手动曲线上判出七个峰，以峰序与峰位作线性拟合，得到自动前测、手动测量和自动后测的有效激发电势分别为 $(11.87\pm0.24)\ \mathrm V$、$(11.70\pm0.21)\ \mathrm V$ 和 $(11.87\pm0.24)\ \mathrm V$，相对参考值 11.61 V 的误差分别为 2.25\%、0.74\% 和 2.22\%。各不确定度仅为拟合标准不确定度。

电流曲线出现一系列峰谷，峰位与峰序近似呈线性关系，说明电子与氩原子的非弹性碰撞存在确定的能量交换尺度，与原子能级离散的解释相符。同时，不能把第一峰电压直接当作激发电势，不能因相关系数接近 1 就忽略峰位判定和系统误差。

额外增加第二次自动测量，是为了验证温度漂移导致了电流测量误差。结果确实有助于验证误差猜想，这也提示我们实验应在充分预热后测量，但也需要在尽可能短的时间内完成手动测量，来避免温度上升导致的温度漂移。

\subsection*{2. 思考题}

\subsubsection*{2.1 第一激发电位的物理含义是什么？}

第一激发电位是以电势差表示原子发生最低激发所需的能量尺度。在理想情况下，电子经过电势差 $U_0$ 加速后获得动能 $eU_0$；当它与基态原子发生非弹性碰撞时，可以把这部分能量传给原子，使原子跃迁到可被激发的最低能态，因此有 $\Delta E=eU_0$。它反映原子内部能级差，而不是电子被电离所需的电离能。

\subsubsection*{2.2 温度对 F-H 管的电流—加速电压曲线有什么影响？}

温度变化会影响阴极的热电子发射：发射电子的数量和初始能量分布改变后，电流基线、峰谷高度及峰宽也会变化。温度若在测量期间持续漂移，前后不同采样点便对应不同工作状态，可能使手动曲线的形状或峰位判定产生偏差。灯丝发射过强还可能使放大器饱和；发射不足则峰谷不明显。

\subsubsection*{2.3 第一个电流峰对应的加速电压为什么不能直接作为氩原子的第一激发电势？}

仪器显示的加速电压不全部对应电子用于激发原子的净动能。不同电极材料引起的接触电势、电子的初始能量、空间电荷和拒斥电场筛选都会影响第一峰的位置；而电流峰值是电子群体在特定碰撞和收集条件下形成的极大值，不是一个理想电子刚达到激发阈值的时刻。
而且本次第一峰在约 27—28 V，而拟合激发电势约为 11.7—11.9 V，两者明显不同。使用 $U_n=A+n\Delta U$ 拟合才能测得较准确的激发电势。
